% !TeX root = part3.tex

\documentclass[../final.tex]{subfiles}

\begin{document}

% ============================================================
% Практическое занятие 3 от 15.09.2026
% ============================================================

\practice{3}{15.09.2026}{Классические и квантовые модели атома}

% ============================================================
% Задача 8. Радиолиния межзвёздного водорода
% ============================================================

\task[8]{Радиолиния межзвёздного атомарного водорода}


% ------------------------------------------------------------
% Условие
% ------------------------------------------------------------

\condition
Найти длину волны излучения межзвёздного атомарного водорода
в радиодиапазоне. Межзвёздный водород находится в основном
электронном состоянии, и его излучение обусловлено переориентацией
спина электрона.


% ------------------------------------------------------------
% Решение
% ------------------------------------------------------------

\solution

Основное состояние атома водорода --- это $1s$-состояние:
%
\begin{equation*}
    n=1,
    \qquad
    l=0.
\end{equation*}

Поэтому орбитальный момент электрона равен нулю, а тонкое
спин-орбитальное расщепление отсутствует. Однако остаётся магнитное
взаимодействие спина электрона со спином протона. Именно оно создаёт
\rconcept{сверхтонкую структуру} основного состояния.


% ------------------------------------------------------------
% Полный спин системы
% ------------------------------------------------------------

\subsection*{Сложение спинов электрона и протона}

Спины электрона и протона равны
%
\begin{equation*}
    S=\frac12,
    \qquad
    I=\frac12.
\end{equation*}

Полный момент
%
\begin{equation*}
    \mathbf F
    =
    \mathbf S
    +
    \mathbf I
\end{equation*}
%
может принимать значения
%
\begin{equation*}
    F=0,
    \qquad
    F=1.
\end{equation*}

Скалярное произведение спинов выражается через $F$:
%
\begin{equation*}
    \mathbf S\cdot\mathbf I
    =
    \frac12
    \left(
        F(F+1)
        -
        S(S+1)
        -
        I(I+1)
    \right)\hbar^2.
\end{equation*}

Для $F=0$:
%
\begin{align*}
    \frac{\langle\mathbf S\cdot\mathbf I\rangle_{F=0}}{\hbar^2}
    &=
    \frac12
    \left(
        0
        -
        \frac34
        -
        \frac34
    \right)
    \\
    &=
    -\frac34.
\end{align*}

Для $F=1$:
%
\begin{align*}
    \frac{\langle\mathbf S\cdot\mathbf I\rangle_{F=1}}{\hbar^2}
    &=
    \frac12
    \left(
        1\cdot2
        -
        \frac34
        -
        \frac34
    \right)
    \\
    &=
    \frac14.
\end{align*}

Следовательно, при переходе между двумя компонентами
сверхтонкой структуры разность спиновых множителей равна
%
\begin{equation*}
    \frac{
        \langle\mathbf S\cdot\mathbf I\rangle_{F=1}
        -
        \langle\mathbf S\cdot\mathbf I\rangle_{F=0}
    }{\hbar^2}
    =
    1.
\end{equation*}


% ------------------------------------------------------------
% Контактное взаимодействие Ферми
% ------------------------------------------------------------

\begin{rbox}[left=12pt,right=12pt,top=12pt,bottom=10pt]{[]}
    \centering
    \begin{tikzpicture}[>=Latex,draw=rtext,every node/.style={text=rtext},font=\small]
        \draw[->] (-1,0) -- (-1,3.2) node[above] {$E$};
        \draw[thick,rconceptcolor] (0,2.6) -- (5,2.6);
        \draw[thick,rresultcolor] (0,0.5) -- (5,0.5);
        \node[right,align=left] at (5.1,2.6) {$F=1$ (триплет)\\$E_1=A/4$};
        \node[right,align=left] at (5.1,0.5) {$F=0$ (синглет)\\$E_0=-3A/4$};
        \draw[->,very thick,rresultcolor] (2,2.6) -- (2,0.5)
            node[midway,left] {$h\nu$};
        \draw[<->] (4.5,0.55) -- (4.5,2.55)
            node[midway,left] {$\Delta E=A$};
        \node[align=center] at (2.4,-0.3) {$1s:\ L=0,\ S=I=\tfrac12$\\$\lambda\approx21\,\text{см}$};
    \end{tikzpicture}
    \captionof{figure}{Сверхтонкое расщепление основного состояния
    водорода и излучательный переход $F=1\to F=0$.
    Здесь $A$ --- положительная константа сверхтонкого взаимодействия;
    энергия отсчитывается от центра тяжести мультиплета.}
    \label{fig:hydrogen-21cm}
\end{rbox}

\subsection*{Контактное взаимодействие Ферми}

Для $s$-состояния электронная плотность в точке расположения протона
не равна нулю. Поэтому основной вклад в сверхтонкое расщепление
даёт контактное взаимодействие Ферми.

В системе \rassumption{СГС} его можно записать в виде
%
\begin{equation*}
    \hat H_{\mathrm{F}}
    =
    \frac{8\pi}{3}
    g_e g_p\mu_B\mu_N
    |\psi_{1s}(0)|^2
    \frac{\hat{\mathbf S}\cdot\hat{\mathbf I}}{\hbar^2}.
\end{equation*}

Здесь
%
\begin{equation*}
    g_e\approx2.0023,
    \qquad
    g_p\approx5.5857,
\end{equation*}
%
$\mu_B$ --- магнетон Бора, а $\mu_N$ --- ядерный магнетон.

Для основного состояния водорода
%
\begin{equation*}
    \psi_{1s}(r)
    =
    \frac{1}{\sqrt{\pi a_0^3}}
    e^{-r/a_0},
\end{equation*}
%
поэтому
%
\begin{equation*}
    |\psi_{1s}(0)|^2
    =
    \frac{1}{\pi a_0^3}.
\end{equation*}

Из предыдущего раздела разность спиновых множителей между $F=1$
и $F=0$ равна единице. Поэтому величина сверхтонкого расщепления
равна
%
\begin{equation*}
    \Delta E_{\mathrm{hfs}}
    =
    \frac{8\pi}{3}
    g_e g_p\mu_B\mu_N
    \frac{1}{\pi a_0^3}.
\end{equation*}

Используем
%
\begin{equation*}
    \mu_B
    =
    9.274\cdot10^{-21}\,
    \frac{\text{эрг}}{\text{Гс}},
    \qquad
    \mu_N
    =
    5.051\cdot10^{-24}\,
    \frac{\text{эрг}}{\text{Гс}},
\end{equation*}
%
и
%
\begin{equation*}
    a_0
    =
    5.292\cdot10^{-9}\,\text{см}.
\end{equation*}

Тогда
%
\begin{equation*}
    \Delta E_{\mathrm{hfs}}
    \approx
    9.4\cdot10^{-18}\,\text{эрг}
    \approx
    5.9\cdot10^{-6}\,\text{эВ}.
\end{equation*}

По порядку величины это совпадает с записанной на семинаре оценкой
%
\begin{equation*}
    \Delta E_{\mathrm{hfs}}
    \sim
    5.7\cdot10^{-6}\,\text{эВ}.
\end{equation*}


% ------------------------------------------------------------
% Длина волны
% ------------------------------------------------------------

\subsection*{Длина волны излучения}

Энергия испускаемого фотона равна
%
\begin{equation*}
    h\nu
    =
    \Delta E_{\mathrm{hfs}}.
\end{equation*}

Следовательно,
%
\begin{equation*}
    \lambda
    =
    \frac{c}{\nu}
    =
    \frac{hc}{\Delta E_{\mathrm{hfs}}}.
\end{equation*}

Подставляя полученное расщепление,
%
\begin{equation*}
    \lambda
    \approx
    21\,\text{см}.
\end{equation*}

Соответствующая частота имеет порядок
%
\begin{equation*}
    \nu
    \approx
    1.42\,\text{ГГц}.
\end{equation*}

Это знаменитая \rconcept{радиолиния нейтрального водорода $21$ см}.


% ------------------------------------------------------------
% Ответ
% ------------------------------------------------------------

\answer

Переход между компонентами сверхтонкой структуры основного состояния
атомарного водорода даёт излучение с длиной волны
%
\begin{equation*}
    \lambda
    \approx
    21\,\text{см}.
\end{equation*}

Искомое излучение действительно лежит в \rresult{радиодиапазоне}.


% ============================================================
% Задача 9. Спин-спиновое расщепление 2p-состояния позитрония
% ============================================================

\task[9]{Спин-спиновое расщепление
\texorpdfstring{$2p$}{2p}-состояния позитрония}


% ------------------------------------------------------------
% Условие
% ------------------------------------------------------------

\condition
Оценить в электронвольтах расщепление $2p$-состояния позитрония,
вызванное взаимодействием спиновых магнитных моментов позитрона
и электрона.


% ------------------------------------------------------------
% Решение
% ------------------------------------------------------------

\solution

Позитроний --- связанное состояние электрона и позитрона.
Обе частицы имеют одинаковую массу $m_e$, поэтому приведённая масса
относительного движения равна
%
\begin{equation*}
    \mu
    =
    \frac{m_em_e}{m_e+m_e}
    =
    \frac{m_e}{2}.
\end{equation*}

Следовательно, боровский радиус позитрония в два раза больше
боровского радиуса атома водорода:
%
\begin{equation*}
    a_{\mathrm{Ps}}
    =
    \frac{m_e}{\mu}a_0
    =
    2a_0.
\end{equation*}


% ------------------------------------------------------------
% Характерный радиус 2p-состояния
% ------------------------------------------------------------

\subsection*{Характерный размер \texorpdfstring{$2p$}{2p}-состояния}

Для водородоподобного $2p$-состояния радиальная функция имеет вид
%
\begin{equation*}
    R_{21}(r)
    \propto
    r\exp\left(-\frac{r}{2a}\right),
\end{equation*}
%
где в случае позитрония $a=a_{\mathrm{Ps}}$.

Радиальная плотность вероятности пропорциональна
%
\begin{equation*}
    w(r)
    \propto
    r^2|R_{21}(r)|^2
    \propto
    r^4e^{-r/a_{\mathrm{Ps}}}.
\end{equation*}

Наиболее вероятный радиус найдём из условия максимума:
%
\begin{equation*}
    \frac{d}{dr}
    \left(
        r^4e^{-r/a_{\mathrm{Ps}}}
    \right)
    =
    0.
\end{equation*}

После вынесения общего положительного множителя получаем
%
\begin{equation*}
    4r^3
    -
    \frac{r^4}{a_{\mathrm{Ps}}}
    =
    0,
\end{equation*}
%
то есть
%
\begin{equation*}
    r_{\mathrm{mp}}
    =
    4a_{\mathrm{Ps}}
    =
    8a_0.
\end{equation*}

Именно этот масштаб расстояния используем для оценки магнитного
взаимодействия.


% ------------------------------------------------------------
% Магнитное диполь-дипольное взаимодействие
% ------------------------------------------------------------

\begin{rbox}[left=12pt,right=12pt,top=12pt,bottom=10pt]{[]}
    \centering
    \begin{tikzpicture}[>=Latex,draw=rtext,every node/.style={text=rtext},font=\small]
        \coordinate (e) at (-2.7,0);
        \coordinate (p) at (2.7,0);
        \draw[dashed] (e) -- (p);
        \fill[rconceptcolor] (e) circle (3pt);
        \fill[rresultcolor] (p) circle (3pt);
        \node[below=6pt] at (e) {$e^-$};
        \node[below=6pt] at (p) {$e^+$};
        \draw[->,thick] (-2.7,-0.8) -- (2.7,-0.8)
            node[midway,below] {$\mathbf r,\quad \mathbf n=\mathbf r/r$};
        \draw[->,thick,rconceptcolor] (e) -- ++(0.6,1.3)
            node[above] {$\mathbf S_-$};
        \draw[->,thick,rresultcolor] (p) -- ++(-0.4,1.3)
            node[above] {$\mathbf S_+$};
        \node[align=center] at (0,2.1)
            {$\boldsymbol\mu_-=-g\mu_B\mathbf S_-/\hbar$\qquad
             $\boldsymbol\mu_+=+g\mu_B\mathbf S_+/\hbar$};
        \node at (0,-1.65) {$a_{\mathrm{Ps}}=2a_0,\qquad r_{\mathrm{mp}}=8a_0$};
    \end{tikzpicture}
    \captionof{figure}{Геометрия спин-спинового взаимодействия в позитронии.
    Направления спинов показаны условно: это схема обозначений оператора,
    а не классические траектории частиц. Взаимодействие зависит
    от ориентации обоих магнитных моментов относительно $\mathbf r$.}
    \label{fig:positronium-spins}
\end{rbox}

\subsection*{Энергия взаимодействия магнитных моментов}

Магнитное диполь-дипольное взаимодействие двух частиц имеет вид
%
\begin{equation*}
    \hat V_{ss}
    =
    \frac{\mu_0}{4\pi}
    \left[
        \frac{\boldsymbol\mu_1\cdot\boldsymbol\mu_2}{r^3}
        -
        3
        \frac{
            (\boldsymbol\mu_1\cdot\mathbf r)
            (\boldsymbol\mu_2\cdot\mathbf r)
        }{r^5}
    \right].
\end{equation*}

Для электрона и позитрона
%
\begin{equation*}
    |\boldsymbol\mu_1|
    \sim
    |\boldsymbol\mu_2|
    \sim
    \mu_B,
\end{equation*}
%
поскольку для обеих частиц $g\approx2$.

В задаче требуется не точное положение всех $2\,{}^3P_J$-компонент,
а только оценка характерного расщепления. Поэтому угловой множитель
в дипольном поле считаем величиной порядка единицы. На оси магнитного
диполя поле имеет порядок
%
\begin{equation*}
    B
    \sim
    \frac{2\mu_B}{r^3},
\end{equation*}
%
а энергия второго магнитного момента в этом поле ---
%
\begin{equation*}
    \Delta E
    \sim
    \mu_B B
    \sim
    \frac{2\mu_B^2}{r^3}.
\end{equation*}

Подставляя характерный радиус $r\sim8a_0$, получаем
%
\begin{equation*}
    \Delta E
    \sim
    \frac{2\mu_B^2}{(8a_0)^3}
    =
    \frac{\mu_B^2}{256a_0^3}.
\end{equation*}


% ------------------------------------------------------------
% Переход к постоянной тонкой структуры
% ------------------------------------------------------------

\subsection*{Оценка через \texorpdfstring{$\alpha$}{alpha}}

Удобнее всего завершить оценку в системе \rassumption{СГС}, где
%
\begin{equation*}
    \mu_B
    =
    \frac{e\hbar}{2m_ec},
    \qquad
    a_0
    =
    \frac{\hbar^2}{m_ee^2},
    \qquad
    \alpha
    =
    \frac{e^2}{\hbar c}.
\end{equation*}

Сначала вычислим комбинацию $\mu_B^2/a_0^3$:
%
\begin{align*}
    \frac{\mu_B^2}{a_0^3}
    &=
    \frac{e^2\hbar^2}{4m_e^2c^2}
    \left(
        \frac{m_ee^2}{\hbar^2}
    \right)^3
    \\
    &=
    \frac{m_e e^8}{4\hbar^4c^2}.
\end{align*}

Из
%
\begin{equation*}
    e^2
    =
    \alpha\hbar c
\end{equation*}
%
следует
%
\begin{equation*}
    e^8
    =
    \alpha^4\hbar^4c^4.
\end{equation*}

Поэтому
%
\begin{equation*}
    \frac{\mu_B^2}{a_0^3}
    =
    \frac14
    \alpha^4m_ec^2.
\end{equation*}

Окончательно
%
\begin{equation*}
    \Delta E
    \sim
    \frac{1}{256}
    \frac14
    \alpha^4m_ec^2
    =
    \frac{\alpha^4m_ec^2}{1024}.
\end{equation*}

Подставляя
%
\begin{equation*}
    \alpha
    \approx
    \frac{1}{137.036},
    \qquad
    m_ec^2
    \approx
    5.11\cdot10^5\,\text{эВ},
\end{equation*}
%
получаем
%
\begin{align*}
    \Delta E
    &\sim
    \frac{1}{1024}
    \left(
        \frac{1}{137.036}
    \right)^4
    5.11\cdot10^5\,\text{эВ}
    \\
    &\approx
    1.4\cdot10^{-6}\,\text{эВ}.
\end{align*}

\remark
Здесь получена именно оценка масштаба спин-спинового расщепления.
Точное вычисление требует усреднения тензорного оператора
$\hat V_{ss}$ по угловой части $2p$-волновой функции и различает
отдельные состояния $2\,{}^3P_J$.


% ------------------------------------------------------------
% Ответ
% ------------------------------------------------------------

\answer

Характерная величина расщепления $2p$-состояния позитрония,
обусловленного взаимодействием спиновых магнитных моментов
электрона и позитрона, составляет
%
\begin{equation*}
    \Delta E
    \sim
    1.4\cdot10^{-6}\,\text{эВ}.
\end{equation*}

То есть эффект имеет масштаб порядка \rresult{одного микроэлектронвольта}.

\end{document}
